% Stable task headings and result-number mapping for the parent master.
% No numbered theorem below is hypothetical.  The three requested tasks
% are reported as incomplete rather than supplied with conditional claims.
\subsection*{C.1. The all-depth ordinary-pair assertion --- INCOMPLETE}
\label{task:ghost-c1}
The asserted necessity of finite support has not been proved or
refuted. Finite-support reconstruction is proved in
Proposition~\ref{prop:ghost-finite-support}. The exact remaining
assertion is that, for every ordinary infinite-support pair,
the continuous section fibre of the full pronilpotent arithmetic
quotient is empty. Theorem~\ref{thm:ghost-fixed-pair-compatible}
now proves that this is equivalent to the existence, for every
such pair, of some integral nilpotency class at which its
global continuous section fibre is empty. This equivalence
keeps the pair fixed and is unconditional.

Any pair contradicting the asserted necessity would have
infinitely ramified, non-point, noncuspidal abelian data in
every coefficient-prime projection, by
Proposition~\ref{prop:ghost-infinite-ramification}. The
known fixed-ramification pro-\(p\) constructions inspected in
this audit do not supply these abelian data; the actual
projection to the \(p\)-primary valuation vector proves
that retaining a rational point at even one coefficient prime
forces finite support. The unramified group
\(H^2(G_{\mathbb Q,\{691,\infty\}},\mathbb Z_{691}(12))\)
injects as a nonzero locally trivial subgroup of absolute
continuous cohomology, by
Proposition~\ref{prop:ghost-h2-691}. Thus automatic vanishing
of arbitrary higher-weight local obstruction classes is not
available; no assertion identifies that subgroup with the
obstruction of an ordinary pair.
Proposition~\ref{prop:ghost-uncountable-kernel} further proves
that this difficulty already occurs in weight two at the
regular coefficient prime three: the integral localization
kernel contains \((\prod\mathbb Z_3)/c_0(\mathbb Z_3)\),
although its fixed-ramification integral group is zero.
For all odd coefficient primes the exact kernel is the
derived inverse limit of absolute finite-coefficient
first cohomology, with the coefficient maps retained.
Proposition~\ref{prop:ghost-ordinary-cup} independently
proves that the continuous cup product of every ordinary
completed pair is zero. Its proof chooses bounded integral
local primitives, then uses finite global primitive fibres
with prescribed inertia restrictions. This verifies the
ordinary-cup conclusion needed from A.27; the unresolved
all-depth question concerns the higher lifting obstructions.
Corollary~\ref{cor:ghost-class2-lift} gives the exact Hall-coordinate
section in the actual arithmetic quotient of class two. Thus
class three is the first possible unrestricted global exclusion
of any ordinary pair. Theorem~\ref{thm:ghost-strict-localization}
now proves the literal continuous localization exact sequence
at every weight, and Proposition~\ref{prop:ghost-class3-residue}
gives the full integral third obstruction of each specified
class-two section, including its derived-limit contribution.
Proposition~\ref{prop:ghost-local-boundary} further proves
that an all-depth section automatically matches the given
ordinary local point at every nonzero valuation prime away
from the coefficient prime. The integral Magnus conjugacy
argument proving that assertion is in
Lemma~\ref{lem:ghost-magnus-conjugacy}. The places of zero
valuation and the equal residue/coefficient prime retain
their actual unforced local fibres. Proposition~\ref{prop:ghost-h1-weight2}
computes the full global weight-two first cohomology as
\(\mathbb Z/24\); Corollary~\ref{cor:ghost-class2-fibre24}
then proves that each ordinary pair has exactly twenty-four
class-two section classes, and a unique such class in each
coefficient prime at least five. The third obstruction thus
has only this finite global freedom.
Theorem~\ref{thm:ghost-first-integral-class3} constructs an
ordinary infinite-support pair for which all these integral
class-two sections have a nonzero third obstruction at a
specified first support prime. Its first unrestricted
integral failure is exactly class three.
Theorem~\ref{thm:ghost-original-first-integral-class3}
achieves the same first depth with the original factorial
exponents of A.44, by treating the coefficient prime three
and its nonzero global central freedom explicitly. This resolves one
concrete first-depth problem without proving that every
ordinary infinite-support pair is excluded at some depth.
Lemma~\ref{lem:ghost-prefix-five} proves the exact
finite-prefix comparison at coefficient five, including
the integral local primitives at the earlier support
primes. Theorem~\ref{thm:ghost-passing-first-class3}
uses it to construct a further original A.44 family
whose first failure is class three even though every
compact elliptic support test is satisfied and a full
affine-cover abelian section exists. These are selected
families. Theorem~\ref{thm:ghost-dense-class3} extends the
construction after every finite legal original A.44 prefix.
Lemma~\ref{lem:ghost-finite-prefix-determines} proves finite
determinacy of the full class-two cochain class modulo every
power of each coefficient prime at least five, by performing
the old-prime division by \(p-1\) integrally before reducing.
Corollary~\ref{cor:ghost-open-dense-class3} consequently gives
an open dense set of first-class-three exclusions in the
original prefix topology, including its compact-passing
subspace. The possible all-depth survivors lie in its closed
complement with empty interior; emptiness or nonemptiness of
that complement has not been established. The all-depth
assertion for every allowed pair remains unresolved.

\subsection*{C.2. Level-dependent survival and the first exclusions --- INCOMPLETE}
\label{task:ghost-c2}
The literal quantifiers of SECTIONS A.65 and A.72 allow the
pair and section to depend on the selected depth or quotient.
Their statements do not assert one fixed pair at every level.
Theorem~\ref{thm:ghost-fixed-depth-survival} now certifies
the full A.65 fixed-depth assertion. Its missing A.66
arithmetic input is reconstructed in
Lemma~\ref{lem:ghost-variable-exponent-image} and
Theorem~\ref{thm:ghost-general-exponent-arithmetic}, with
the exact local shell images and every radical/CRT
intersection checked. The certification uses the proved
strict localization morphism, a uniform finite unramified
annihilator including coefficient two, and an actual
Bockstein correction retaining the stated finite Magnus
congruence. Thus for every fixed integral class there
is one ordinary infinite-support pair surviving that class.
The stronger assertion that no one pair survives every
integral nilpotency class remains unresolved, exactly as
specified under C.1. The displayed A.65 outputs converge
to the rational pair \((3,-2)\), by
Proposition~\ref{prop:ghost-a65-limit}. A common-cofinal family
of A.72 sections retaining the point-three component at every
finite quotient likewise has abelian class exactly
\((3,-2)\), since the quotients include every abelian exponent.
Neither calculation excludes other choices of deeper sections.

Theorem~\ref{thm:ghost-factorial-selmer32} supplies a fully
specified residue-prescribed variant of A.44 whose first
failure of local matching is exactly class two, detected by
the indicated geometric group of order \(32\). The proof
does not claim failure of unrestricted global lifting.
Proposition~\ref{prop:ghost-1129-unramified} proves that
the prime \(1129\) contributes no local third obstruction
for any global class-two lift of this witness, and gives a
full local pro-\(\ell\) lift there for every \(\ell\ne1129\).
The class-three residue criterion retains the support
primes and their convergence condition as the unresolved
global calculation for this particular pair.
No smallest excluding level has been established for every
output of the canonical least-choice A.44 recursion.

Theorem~\ref{thm:ghost-first-integral-class3} supplies a
different completely specified variant, with exponents
\(5(j+4)!\) and a first-prime Frobenius condition in the
point-three Heisenberg field. Its first unrestricted integral
failure is exactly class three: every integral class-two lift
has a nonzero coefficient-prime-five local third obstruction
at that first support prime. The reduction of each such
obstruction modulo twenty-five is nonzero. The theorem keeps
the integral domain of these reductions; arbitrary finite
coefficient class-two sections can have additional ramified
classes, which this calculation does not classify.
Theorem~\ref{thm:ghost-original-first-integral-class3}
strengthens this construction to a choice of primes allowed
by the original A.44 recursion, with exponents \((j+4)!\).
Its obstruction occurs at coefficient prime three, modulo
nine, and the first-prime rule is an explicit cubic-radical
Frobenius test. The question for every allowed A.44 choice,
including an independently prescribed least-choice output,
remains unresolved.

Theorem~\ref{thm:ghost-passing-first-class3} settles
the unrestricted third obstruction for a selected family
inside the compact-passing domain of
Proposition~\ref{prop:ghost-elliptic-passing}. It keeps
all original factorial exponents, freezes the old
integral primitive values using
Lemma~\ref{lem:ghost-prefix-five}, and imposes a
compatible nonzero central Heisenberg Frobenius at
the second positive support prime \(q_1\).
All twenty-four integral class-two lifts fail there
at coefficient five, already modulo twenty-five.
Every compact elliptic test and the full affine
abelian cover quotient still have the stated surviving
classes. The unrestricted calculation for all the
other compact-passing outputs is not settled.
Theorem~\ref{thm:ghost-dense-class3} nevertheless places such
an exclusion after every finite legal original prefix,
and Corollary~\ref{cor:ghost-open-dense-class3} strengthens
this to a finite extension whose every legal continuation
is excluded. This is an open dense statement in the exact
original-prefix topology, with the factorial exponents
unchanged. It does not decide the remaining closed domain.

SECTIONS A.137 and A.139 impose the specified local matching
conditions and obtain finite Brauer reciprocity exclusions;
their omitted A.134 arithmetic construction prevents a full
independent certification from this attachment alone. They
refer to the same harmonic pair: A.137 names a detector of
class seventeen, and A.139 one of class eight. The latter's
actual group \((\mathbb Z/2)^8\rtimes\mathbb Z/16\), its
order \(4096\), and exact class eight have been independently
verified by the nine Schreier generators and the powers of
the eight-cycle's augmentation operator. Detector class alone
does not establish the smallest failing class. A.207
and A.209 assert a global class-four lift and failure of local
matching at class two; forgetting the local conjugacy
conditions is the exact map relating those assertions. The
order-32 local-matching obstruction has been independently
proved here. A.196 is over \(\mathbb Q(\sqrt[3]3)\), and
its local class modulo ninth powers realizes no ordinary
point or cusp, by Proposition~\ref{prop:ghost-a196-scope}.
Its finite supported obstruction and its claimed continuous
absolute obstruction must retain their respective source
Galois groups. The independent weight-twelve example of
Corollary~\ref{cor:ghost-integral-extension} proves that
absolute finite coefficient splitting can indeed coexist with
continuous nonsplitting over \(\mathbb Q\), without asserting
that this example arises from a unit-equation pair.

\subsection*{C.3. A uniform rigid target --- INCOMPLETE}
\label{task:ghost-c3}
Proposition~\ref{prop:ghost-elliptic-target} identifies an
actual degree-\(27\) finite cover excluding the family of
Theorem~\ref{thm:ghost-original-first-integral-class3}:
it is the elliptic curve \(y^2=x^3+9\) with twenty-seven
geometric punctures, with its explicit cover map through
\(y^2=x^3-243\) and \(3u^3-2v^3=1\). All three
mod-three central classes are accounted for: two already
fail the third obstruction at two, and every surviving
full section must factor through the single pointed cover.
At the constructed first support prime its ordinary point
has no rational lift to that cover. This proves the full
pro-three section exclusion using the actual cover map.
The compactification kills boundary inertia that maps
nontrivially to \(U\); the proof therefore retains the
open curve and does not claim exclusion from its compact
elliptic quotient alone.

Propositions~\ref{prop:ghost-elliptic-selmer}--\ref{prop:ghost-elliptic-support-test}
compute the exact three-isogeny Selmer groups and the
absolute continuous cohomology
\(H^1(G_{\mathbb Q},T_3E')=\mathbb Z_3\kappa(6,15)
\oplus(\mathbb Z/3)\kappa(0,3)\).
The exact translated isogeny identity then forces, for
every original A.44 output admitting a full pro-three
section, that \((-2,1)\) have prime-to-three order
at every positive support prime \(q_j\), and that
\((3,-6)\) have prime-to-three order at every negative
support prime \(r_j\). Proposition~\ref{prop:ghost-elliptic-passing}
constructs an original A.44 family satisfying every
one of these finite elliptic tests. The remaining case
for this compact elliptic calculation is therefore
explicitly nonempty; existence or exclusion of a full
section for it remains unresolved.

Lemma~\ref{lem:ghost-elliptic-two-fibre} proves that the
compact comparison at two adds no restriction on the two
global coefficients. Proposition~\ref{prop:ghost-genjac-image}
identifies the full affine-cover geometric abelianization
with the Tate module of the generalized Jacobian, retains
its exact map to the two original Kummer coordinates, and
computes its image by the norms from the three boundary
algebras. Proposition~\ref{prop:ghost-boundary-hasse-norm}
proves that these are degree-nine fields with normal-closure
group \(H_3\rtimes C_2\), and that their norm-one tori satisfy
both the Hasse norm principle and weak approximation.
Proposition~\ref{prop:ghost-original-genjac-survival} then
constructs an original, unscaled A.44 output with an actual
full profinite section of
\(\pi_1(U)/\overline{[K,K]}\), where
\(K=\pi_1(Y_{\overline{\mathbb Q}})\), factoring through the
generalized-Jacobian quotient of \(Y\). Its proper class is
zero. The remaining cover-lifting kernel is the explicitly
retained \(\overline{[K,K]}\).
Proposition~\ref{prop:ghost-all-passing-genjac} strengthens
this survival statement to every original A.44 output
satisfying the compact support-order tests: the boundary
norms are surjective over \(\mathbb Q_3\), so the extra
congruence used in the first construction is unnecessary.
Theorem~\ref{thm:ghost-passing-first-class3} and
Corollary~\ref{cor:ghost-metabelian-cover-exclusion}
then exclude a selected subfamily by the actual
metabelian quotient of \(Y\). Its abelian cover
section fibre is nonempty, but every member has empty
lift fibre through
\(\pi_1(Y)/\overline{[[K,K],[K,K]]}\).
The exact map to the downstairs class-three quotient
comes from
\(\overline{[[K,K],[K,K]]}\subset\Gamma_4\Pi\).
Thus the retained commutator problem is genuinely
nonempty and can obstruct, while a universal exclusion
or full lift for all the remaining compact-passing
outputs is still unresolved.

The order-32 quotient above excludes the residue-prescribed
family. The complete attached A.72 argument has been checked:
for every one fixed finite geometric quotient \(Q\), it
constructs an ordinary infinite-support pair with a locally
matching section in
\[
 E/(N\cap[\Pi,\Pi])
       =E_Q\mathbin{\times}_{E_{Q^{\rm ab}}}E_{\rm ab},
       \qquad N=\ker(\Pi\to Q),
\]
whose \(Q\)-component is the section of three. The equality
is induced by the diagonal quotient map; its kernel is the
displayed intersection, and compatibility over \(Q^{\rm ab}\)
allows any pair of lifts to be adjusted by the two kernels.
Therefore no single finite geometric quotient uniformly
excludes all these A.72 ordinary ghosts. A finite collection
of quotient tests has the same limitation, by taking its
common finite quotient.

A geometrically connected finite \emph{etale} cover \(Y\to U\) of degree
\(d\) gives an open subgroup of the full fundamental group.
Its geometric subgroup has index \(d\) in the free profinite
group of rank two and therefore has rank \(d+1\), by the
Schreier generators and relations, so is infinite. A.72 constructs matching sections of the
specified finite-quotient fibre product; it does not construct
a section of \(\pi_1(Y)\) or solve every quotient test of
that infinite group. No proof has been obtained of a uniform
exclusion through a full finite-cover fundamental group, a
configuration-space target, or a modular target, and no
proof has been obtained that every target of those specified
types fails. This is the precise remaining target problem.

\paragraph{Lead 3 --- partially proved, main assertion INCOMPLETE.}
The finite-support reconstruction and the equivalence between
fixed-pair integral survival at every depth and actual
pronilpotent survival are proved. The asserted implication
from this survival to finite support remains unresolved.

\paragraph{Lead 6 --- section component INCOMPLETE.}
The finite quotient of order \(32\) and its cover-to-Jacobian
map give a concrete geometric target for one specified
family. The fixed elliptic cover and its metabelian quotient
give further exact exclusions, including a family that
survives the whole affine geometric abelianization of that
cover. A.72 proves the precisely stated failure of any
single finite geometric quotient as a uniform test. The
comparison to the full infinite fundamental group of a
finite cover or to the other requested rigid targets remains
unresolved. The GT component of this lead is audited in
Part A and is not replaced by a section-theoretic assertion.

% Required parent-master numbering map:
% C.1.1 prop:ghost-finite-support
% C.1.2 prop:ghost-infinite-ramification
% C.1.3 lem:ghost-higher-z1
% C.1.4 lem:ghost-hall-model
% C.1.5 lem:ghost-primitive-family
% C.1.6 prop:ghost-compact-fibres
% C.1.7 thm:ghost-fixed-pair-compatible
% C.1.8 prop:ghost-h2-691
% C.1.9 cor:ghost-integral-extension
% C.1.10 prop:ghost-uncountable-kernel
% C.1.11 prop:ghost-ordinary-cup
% C.1.12 cor:ghost-class2-lift
% C.1.13 thm:ghost-strict-localization
% C.1.14 prop:ghost-tame-cochain
% C.1.15 prop:ghost-class3-residue
% C.1.16 lem:ghost-magnus-conjugacy
% C.1.17 prop:ghost-local-boundary
% C.1.18 prop:ghost-h1-weight2
% C.1.19 cor:ghost-class2-fibre24
% C.1.20 lem:ghost-prefix-five
% C.1.21 lem:ghost-finite-prefix-determines
% C.2.1 prop:ghost-a65-limit
% C.2.2 thm:ghost-factorial-selmer32
% C.2.3 prop:ghost-a196-scope
% C.2.4 prop:ghost-1129-unramified
% C.2.5 thm:ghost-first-integral-class3
% C.2.6 thm:ghost-original-first-integral-class3
% C.2.7 thm:ghost-passing-first-class3
% C.2.8 thm:ghost-dense-class3
% C.2.9 cor:ghost-open-dense-class3
% C.2.10 lem:ghost-variable-exponent-image
% C.2.11 thm:ghost-general-exponent-arithmetic
% C.2.12 thm:ghost-fixed-depth-survival
% C.3.1 prop:ghost-elliptic-target
% C.3.2 prop:ghost-elliptic-selmer
% C.3.3 prop:ghost-elliptic-h1
% C.3.4 prop:ghost-elliptic-support-test
% C.3.5 prop:ghost-elliptic-passing
% C.3.6 lem:ghost-elliptic-two-fibre
% C.3.7 prop:ghost-genjac-image
% C.3.8 prop:ghost-boundary-hasse-norm
% C.3.9 prop:ghost-original-genjac-survival
% C.3.10 prop:ghost-all-passing-genjac
% C.3.11 cor:ghost-metabelian-cover-exclusion
